\documentclass[UTF8,12pt]{ctexbook}

% 支持从项目根目录或当前笔记目录编译。
\makeatletter
\def\input@path{{../}{../../}{../../../}{../../../../}}
\makeatother
\usepackage{researchnotes}

\title{概率论学习思考与习题}
\author{}
\date{}

\begin{document}
\maketitle

\frontmatter
\tableofcontents

\chapter{绪论}

概率论与统计学都研究随机现象和不确定性，但侧重点不同。\keyterm{概率论（probability theory）}从给定的假设和规则出发，研究各种结果出现的可能性及其规律；\keyterm{统计学（statistics）}研究如何收集和利用数据，描述已经观察到的现象、判断未知情况，并评价判断的可靠程度。\impo{前者侧重从模型推导可能的结果，后者的一个主要任务是利用数据认识模型及其未知部分。}

以重复抛掷硬币为例，假定每次正面出现的机会保持不变，各次抛掷的结果互不影响。如果还知道正反面出现的机会相同，研究“抛掷 100 次，恰好出现 60 次正面的可能性有多大”，就是典型的概率论问题。如果正面出现的机会未知，实际抛掷 100 次，记录到 60 次正面，再据此估计这一机会，并考察估计可能有多大误差，就是典型的统计学问题。

\impo{“记录中正面占六成”是在描述已观察到的事实，并不意味着硬币每次正面出现的真实机会一定是六成。}即使正反面的机会相同，抛掷 100 次也不保证正面恰好出现 50 次。因此，根据数据判断硬币的性质时，需要考虑偶然波动，不能仅凭次数偏离 50 次就断言硬币不公平。

概率论为统计方法提供数学基础，使我们能够研究数据和判断中的随机波动；统计学则利用数据认识未知情况，并检查模型是否适合实际问题。\keyterm{数理统计（mathematical statistics）}着重研究统计方法的数学原理及其性质，是统计学的重要组成部分。两者合在一起学习，体现了从认识随机规律到利用数据作出判断的知识联系。\impo{学习时应分清模型的假设、实际的观察记录，以及据此能够得到的结论。}

本笔记以缪柏其、张伟平主编的《概率论与数理统计》（高等教育出版社，2022 年第 1 版）为学习线索，按教材章节顺序整理学习思考与习题解答。

\mainmatter

\chapter{事件及其概率}
\label{chap:study-events}

\section{概率论简史}
\label{sec:study-probability-history}

\subsection{中止比赛的赌本公平分配}
\label{subsec:points-fair-allocation}

甲、乙各拿出 $150$ 元，共 $300$ 元作为赌本，约定先赢三局者拿走全部赌本。每局没有平局，双方获胜的机会相同，各局结果相互独立。现在甲已赢两局、乙已赢一局，比赛却因外部原因中止。这 $300$ 元应怎样公平分配？关键在于：当前比分影响双方继续比赛后获胜的机会，不能直接按已赢局数之比 $2:1$ 分配。

\subsubsection{问题设定与分配原则}
\label{subsubsec:points-problem}

先把问题一般化。设 $A,B$ 各投入 $M/2$ 元，总赌本为 $M>0$ 元；比赛采用先赢 $s$ 局者获胜的规则，其中 $s$ 为正整数，胜者取得全部赌本，败者所得为 $0$。每局无平局，各局相互独立，$A$ 每局获胜的概率为已知常数 $p\in(0,1)$，$B$ 为 $1-p$。中止时双方分别已赢 $a,b$ 局，其中 $a,b\in\{0,1,\ldots,s-1\}$，尚未产生最终胜者。假定中止原因不提供关于未来胜负的额外信息，若比赛继续，后续各局仍遵循原有假设。

这里将“公平”规定为\keyterm{保持双方在当前信息下的期望所得}：提前分得的金额，应等于假想比赛继续至结束时各自所得的条件期望（conditional expectation）。因此，需要计算的是双方最终获胜的概率。记
\begin{equation}
  r=s-a,\qquad t=s-b,
  \label{eq:points-remaining}
\end{equation}
分别为 $A,B$ 还需取得的胜局数，并用 $Q(r,t)$ 表示在当前状态下 $A$ 最终获胜的概率。由于胜率已知且固定、各局独立，过去各局的具体先后顺序不影响这一概率。

$A$ 最终所得以概率 $Q(r,t)$ 为 $M$，以概率 $1-Q(r,t)$ 为 $0$，故其期望为 $M Q(r,t)$；$B$ 的期望同理。双方应分得
\begin{equation}
  x_A=M Q(r,t),\qquad x_B=M\bigl[1-Q(r,t)\bigr].
  \label{eq:points-allocation}
\end{equation}
这既保证金额非负，也保证 $x_A+x_B=M$。分配金额包含原始投入，净收益则分别为 $x_A-M/2$ 和 $x_B-M/2$。

\subsubsection{补齐比赛与一般公式}
\label{subsubsec:points-binomial}

为了计算 $Q(r,t)$，令
\begin{equation}
  N=r+t-1.
  \label{eq:points-horizon}
\end{equation}
若还没有胜者，$A$ 至多再赢 $r-1$ 局，$B$ 至多再赢 $t-1$ 局，合计至多进行 $r+t-2=N-1$ 局。因此，比赛最多再进行 $N$ 局便会结束。

现在预先生成 $N$ 个相互独立的未来比赛结果，每个结果中 $A$ 获胜的概率均为 $p$。真实比赛仍在某一方达到所需胜局数时结束，余下结果只用于把序列假想补齐，不改变已经确定的胜者。这样便可在固定的 $N$ 局中计数，而不必分别讨论比赛在不同局数结束的情形。

设随机变量 $K$ 表示这 $N$ 局中 $A$ 的胜局数。恰有 $k$ 局获胜的位置共有 $\binom Nk$ 种选法，每种结果序列的概率均为 $p^k(1-p)^{N-k}$，所以 $K$ 服从\keyterm{二项分布（binomial distribution）}，且
\begin{equation}
  \mathbb P(K=k)=\binom Nk p^k(1-p)^{N-k},
  \qquad k=0,1,\ldots,N.
  \label{eq:points-binomial-mass}
\end{equation}

\begin{proposition}[赌本分配的一般公式]
\label{prop:points-general-shares}
在上述假设下，令 $r=s-a$、$t=s-b$、$N=r+t-1$，则
\begin{equation}
  Q(r,t)=\sum_{k=r}^{N}\binom Nk p^k(1-p)^{N-k},
  \label{eq:points-binomial-tail}
\end{equation}
双方按当前信息下的期望所得应分得
\begin{equation}
  x_A=M\sum_{k=r}^{N}\binom Nk p^k(1-p)^{N-k},
  \qquad x_B=M-x_A.
  \label{eq:points-general-shares}
\end{equation}
\end{proposition}

\begin{proof}
关键是说明：$A$ 在真实比赛中获胜，当且仅当补齐后 $K\ge r$。若 $A$ 获胜，补齐后自然仍至少有 $r$ 个胜局；若 $B$ 获胜，则补齐后 $B$ 至少有 $t$ 个胜局，因而 $A$ 至多赢 $N-t=r-1$ 局。比赛必产生唯一胜者，故上述两个事件等价。于是，对式~\eqref{eq:points-binomial-mass} 中 $k=r,\ldots,N$ 的概率求和，再乘以赌本 $M$，便得到分配公式。
\end{proof}

\subsubsection{计算示例与适用范围}
\label{subsubsec:points-examples}

\begin{example}
\label{ex:points-original}
回到开头的问题：总赌本为 $300$ 元，双方每局胜率相等，先赢三局者获胜，中止时比分为 $2:1$。求双方应分得的金额。
\end{example}

\begin{solve}
此时 $r=1$、$t=2$，最多还需 $N=2$ 局。补齐后的两局结果为 $AA,AB,BA,BB$，每种概率均为 $1/4$；其中字母依次表示各局胜者，只有 $BB$ 对应 $B$ 最终获胜。注意，$AA,AB$ 中比赛其实在第一局后就已结束，第二个字母只是为计数而补上的结果。因此，$Q(1,2)=3/4$，应分给双方
\[
  x_A=300\times\frac34=225\text{ 元},\qquad
  x_B=300\times\frac14=75\text{ 元}.
\]
这也可直接看出：$B$ 必须连赢两局才能获胜。若每局 $A$ 的胜率改为一般的 $p$，则 $B$ 最终获胜的概率为 $(1-p)^2$，相应分配为 $x_A=300[1-(1-p)^2]$、$x_B=300(1-p)^2$。
\end{solve}

\begin{example}
\label{ex:points-biased}
双方各投入 $500$ 元，约定先赢四局者获得全部赌本。各局独立，$A$ 每局获胜的概率为 $3/5$；中止时 $A$ 已赢一局，$B$ 已赢两局。求公平分配金额。
\end{example}

\begin{solve}
此时 $M=1000$、$r=3$、$t=2$、$N=4$。由式~\eqref{eq:points-binomial-tail}，
\[
  Q(3,2)=\binom43\left(\frac35\right)^3\frac25
          +\binom44\left(\frac35\right)^4
        =\frac{297}{625}.
\]
所以 $A$ 分得 $1000\times297/625=475.2$ 元，$B$ 分得 $524.8$ 元。虽然 $A$ 每局胜率较高，但距离最终胜利还差更多局，说明分配同时取决于每局胜率和当前比分。
\end{solve}

上述“公平”以保持期望所得为准，投入相等并不意味着中止时应平分赌本。公式要求各局独立、胜率已知且固定，中止原因也不提供额外信息；若这些条件不成立，就需另建模型，不能把当前胜局比例直接当作每局胜率套用。

\section{随机试验和随机事件}
\label{sec:study-random-experiments}

\subsection{德摩根律：复合事件的对立事件}
\label{subsec:study-de-morgan}

描述随机事件时，经常需要把一个条件整体否定。例如，“两个事件至少发生一个”的否定是什么？“两个事件都发生”的否定又是什么？\impo{\keyterm{德摩根律（De Morgan's laws）}说明：对一个并集或交集取补集时，需要交换并与交，并对其中每个集合分别取补集。}在概率论中，它把复合事件的“不发生”转化为若干较简单事件之间的关系，是寻找对立事件的基本依据。

先明确记号。随机试验所有可能结果组成的集合称为\keyterm{样本空间（sample space）}，记作 $\Omega$；其中的一个结果记作 $\omega$，称为样本点。事件 $A$ 用 $\Omega$ 中满足指定条件的样本点组成的集合表示，实际结果 $\omega\in A$ 就表示事件 $A$ 发生。相对于同一个样本空间，事件 $A$ 的\keyterm{对立事件（complementary event）}是
\[
  A^{\mathrm c}=\Omega\setminus A
  =\{\omega\in\Omega:\omega\notin A\},
\]
也就是“$A$ 不发生”，其集合称为 $A$ 的补集，有时也记作 $\overline A$。对于两个事件 $A,B$，\keyterm{并（union）} $A\cup B$ 表示“$A,B$ 至少有一个发生”，允许二者同时发生；\keyterm{交（intersection）} $A\cap B$ 表示“$A,B$ 同时发生”。下面所有补集都相对于固定的 $\Omega$ 而言。

\begin{proposition}[两个事件的德摩根律]
\label{prop:study-de-morgan}
设 $A,B$ 是同一随机试验中的两个事件，则
\begin{align}
  (A\cup B)^{\mathrm c}
    &=A^{\mathrm c}\cap B^{\mathrm c},
    \label{eq:study-de-morgan-union}\\
  (A\cap B)^{\mathrm c}
    &=A^{\mathrm c}\cup B^{\mathrm c}.
    \label{eq:study-de-morgan-intersection}
\end{align}
\end{proposition}

式~\eqref{eq:study-de-morgan-union} 表示：\keyterm{“至少一个发生”的对立事件是“一个也不发生”}，因此必须让 $A$ 不发生且 $B$ 不发生。式~\eqref{eq:study-de-morgan-intersection} 表示：\keyterm{“都发生”的对立事件是“至少一个不发生”}，只要 $A,B$ 中有一个没有发生，就足以使“都发生”不成立。这是逻辑中否定与“或”“且”的关系：否定一个“或”条件得到两个否定条件的“且”，否定一个“且”条件得到两个否定条件的“或”。

\begin{example}
\label{ex:study-de-morgan-die}
掷一枚六面骰子一次，取 $\Omega=\{1,2,3,4,5,6\}$，样本点表示朝上一面的点数。设 $A$ 为“点数为偶数”，$B$ 为“点数大于 $3$”。分别写出 $A\cup B$、$A\cap B$ 及其对立事件，并核对德摩根律。
\end{example}

\begin{solve}
先按条件列出样本点：
\[
  A=\{2,4,6\},\qquad B=\{4,5,6\},\qquad
  A^{\mathrm c}=\{1,3,5\},\qquad B^{\mathrm c}=\{1,2,3\}.
\]
于是
\[
  \begin{aligned}
    A\cup B&=\{2,4,5,6\},&
    (A\cup B)^{\mathrm c}&=\{1,3\}
      =A^{\mathrm c}\cap B^{\mathrm c},\\
    A\cap B&=\{4,6\},&
    (A\cap B)^{\mathrm c}&=\{1,2,3,5\}
      =A^{\mathrm c}\cup B^{\mathrm c}.
  \end{aligned}
\]
第一行中，“既不是偶数，也不大于 $3$”只有点数 $1,3$；第二行中，“不是偶数，或者不大于 $3$”包括点数 $1,2,3,5$。特别地，点数为 $2$ 时 $A$ 发生而 $B$ 不发生，属于“$A,B$ 不都发生”，却不属于“$A,B$ 都不发生”。\impo{因此，“不都发生”不能理解为“都不发生”，也不能理解为“恰有一个不发生”，因为它还包括二者都不发生的情形。}
\end{solve}

德摩根律同样适用于有限多个事件。设 $n\ge1$ 为整数，$A_1,\ldots,A_n$ 是同一样本空间中的事件，则
\begin{equation}
  \left(\bigcup_{i=1}^{n}A_i\right)^{\mathrm c}
    =\bigcap_{i=1}^{n}A_i^{\mathrm c},
  \qquad
  \left(\bigcap_{i=1}^{n}A_i\right)^{\mathrm c}
    =\bigcup_{i=1}^{n}A_i^{\mathrm c}.
  \label{eq:study-de-morgan-finite}
\end{equation}
第一式说，“至少一个事件发生”的否定是“所有事件都不发生”；第二式说，“所有事件都发生”的否定是“至少一个事件不发生”。从样本点看，前者把“存在一个 $i$ 使 $\omega\in A_i$”的否定改写为“对每个 $i$ 都有 $\omega\notin A_i$”；后者把“对每个 $i$ 都有 $\omega\in A_i$”的否定改写为“存在一个 $i$ 使 $\omega\notin A_i$”。这既证明了有限情形，也说明同样的论证对可数无穷事件列 $A_1,A_2,\ldots$ 仍成立：
\begin{equation}
  \left(\bigcup_{i=1}^{\infty}A_i\right)^{\mathrm c}
    =\bigcap_{i=1}^{\infty}A_i^{\mathrm c},
  \qquad
  \left(\bigcap_{i=1}^{\infty}A_i\right)^{\mathrm c}
    =\bigcup_{i=1}^{\infty}A_i^{\mathrm c}.
  \label{eq:study-de-morgan-countable}
\end{equation}

\impo{德摩根律本身是\keyterm{事件之间的恒等式}，不要求事件互斥或相互独立，也不涉及各事件发生的概率大小。}引入概率后，记 $\mathbb P(E)$ 为事件 $E$ 的概率，结合对立事件的概率关系 $\mathbb P(E^{\mathrm c})=1-\mathbb P(E)$，式~\eqref{eq:study-de-morgan-finite} 给出
\begin{equation}
  \mathbb P\!\left(\bigcup_{i=1}^{n}A_i\right)
    =1-\mathbb P\!\left(\bigcap_{i=1}^{n}A_i^{\mathrm c}\right).
  \label{eq:study-de-morgan-probability}
\end{equation}
\impo{因此，求“至少一个发生”的概率时，可以先求“一个也不发生”的概率，再用 $1$ 减去它。}\impo{不过，交事件的概率不能仅凭德摩根律写成各事件概率的乘积；能否相乘还需要相应的独立性条件。}记忆时可概括为：\keyterm{整体取补，并交互换，逐项取补}；理解时则应回到“至少一个”与“所有”的否定关系。

\section{概率的定义和性质}
\label{sec:study-probability-definition}

\begin{definition}
\label{def:study-classical-probability}
设随机试验的样本空间为 $\Omega=\{\omega_1,\ldots,\omega_n\}$，其中 $n\ge1$，各样本点互不相同。若每个样本点发生的概率均为 $1/n$，则称这一概率模型为\keyterm{古典概型（classical probability model）}。\impo{它具有两个基本特征：样本点有限，且各样本点等可能。}对任意事件 $A\subseteq\Omega$，有
\begin{equation}
  \mathbb P(A)=\frac{|A|}{|\Omega|}
  =\frac{\text{属于事件 }A\text{ 的样本点数}}{\text{全部样本点数}}.
  \label{eq:study-classical-probability}
\end{equation}
其中 $|A|$ 表示集合 $A$ 中的元素个数。
\end{definition}

上述公式来自概率的有限可加性：事件 $A$ 由 $|A|$ 个互斥的单点事件组成，每个单点事件的概率为 $1/n$，故总概率为 $|A|/n$。例如，掷一枚均匀骰子，六种点数等可能，“出现偶数”包含 $2,4,6$ 三个样本点，因此其概率为 $3/6=1/2$。应用古典概型前，应先明确样本空间并核实各样本点的等可能性，不能仅因结果个数有限就直接按数量之比计算概率。

\subsection{生日问题：对立事件与对数近似}
\label{subsec:study-birthday-problem}

% 教材例 1.7：生日问题（birthday problem）。
\begin{example}
\label{ex:study-birthday}
试求 $r$ 人中至少有 $2$ 人生日相同的概率，一年按 $365$ 天计，其中 $r$ 为正整数。假设各人的生日相互独立，且每个人的生日在这 $365$ 天上均匀分布。
\end{example}

\begin{solve}
记 $A=\{\text{至少两人生日相同}\}$，则 $A^{\mathrm c}=\{\text{所有人的生日两两不同}\}$。若 $r\ge366$，必有至少两人生日相同，故 $\mathbb P(A)=1$；以下设 $1\le r\le365$。将日期编号为 $1,\ldots,365$，样本空间为 $\Omega=\{1,\ldots,365\}^{r}$，其中一个样本点 $(x_1,\ldots,x_r)$ 记录每个人的生日。由题设，$365^r$ 个样本点等可能。要使生日两两不同，各人的生日依次有 $365,364,\ldots,365-r+1$ 种选择，故
\begin{equation}
  \mathbb P(A)=1-\frac{365!}{(365-r)!\,365^r}
  =1-\prod_{k=1}^{r-1}\left(1-\frac{k}{365}\right).
  \label{eq:study-birthday-exact}
\end{equation}
其中 $r=1$ 时空乘积取 $1$，因而 $\mathbb P(A)=0$。

\impo{取自然对数可将连乘化为求和，再利用 $x$ 较小时的近似 $\ln(1-x)\approx-x$ 简化估算。}于是
\begin{equation}
  \begin{aligned}
    \ln\mathbb P(A^{\mathrm c})
      &=\sum_{k=1}^{r-1}\ln\left(1-\frac{k}{365}\right)\\
      &\approx-\frac1{365}\sum_{k=1}^{r-1}k
      =-\frac{r(r-1)}{730}.
  \end{aligned}
  \label{eq:study-birthday-log}
\end{equation}
取指数还原，得到
\begin{equation}
  \mathbb P(A)\approx1-\exp\left(-\frac{r(r-1)}{730}\right).
  \label{eq:study-birthday-approx}
\end{equation}
此近似要求各项 $k/365$ 较小且累计误差可以接受。$r=30,50$ 时，近似值分别为 $0.6963,0.9651$；按精确公式~\eqref{eq:study-birthday-exact} 计算，相应结果为 $0.7063,0.9704$（保留四位小数）。
\end{solve}

\subsection{加法定理与容斥原理：重叠事件的概率}
\label{subsec:study-inclusion-exclusion}

求若干事件中“至少有一个发生”的概率，就是求这些事件的并集的概率。若事件之间存在重叠，直接相加各事件的概率，就可能把共同部分重复计算。概率的\keyterm{加法定理（addition theorem）}通过交替减去和补回各重交集的概率，修正这种重复计数；它是\keyterm{容斥原理（inclusion--exclusion principle）}在概率计算中的形式。以下各事件均属于同一概率空间，用 $\mathbb P$ 表示概率；教材中的 $A_iA_j$ 表示交集 $A_i\cap A_j$，不是概率的乘积。

先考察两个事件。记 $B\setminus A=B\cap A^{\mathrm c}$，即“$B$ 发生而 $A$ 不发生”。由于 $A\cup B$ 是 $A$ 与 $B\setminus A$ 的不交并，而 $B$ 是 $B\setminus A$ 与 $A\cap B$ 的不交并，由概率的有限可加性，
\[
  \mathbb P(A\cup B)=\mathbb P(A)+\mathbb P(B\setminus A),
  \qquad
  \mathbb P(B)=\mathbb P(B\setminus A)+\mathbb P(A\cap B).
\]
消去 $\mathbb P(B\setminus A)$，得到
\begin{equation}
  \mathbb P(A\cup B)
    =\mathbb P(A)+\mathbb P(B)-\mathbb P(A\cap B).
  \label{eq:study-inclusion-exclusion-two}
\end{equation}
当 $A,B$ 互斥，即 $A\cap B=\varnothing$ 时，交集项为零，公式退化为互斥事件的加法公式。

对于三个事件，先对 $A\cup B$ 与 $C$ 使用式~\eqref{eq:study-inclusion-exclusion-two}，再由分配律 $(A\cup B)\cap C=(A\cap C)\cup(B\cap C)$ 对交集部分使用同一公式，得
\begin{equation}
  \begin{aligned}
    \mathbb P(A\cup B\cup C)
      ={}&\mathbb P(A)+\mathbb P(B)+\mathbb P(C)\\
         &-\mathbb P(A\cap B)-\mathbb P(A\cap C)-\mathbb P(B\cap C)\\
         &+\mathbb P(A\cap B\cap C).
  \end{aligned}
  \label{eq:study-inclusion-exclusion-three}
\end{equation}
三者交集为何要加回来？属于三个事件共同部分的样本点，在单事件项中计入三次，又在两两交集项中减去三次，因此还要补回一次；只属于其中两个事件的样本点则计入两次、减去一次，已经恰好保留一次。

\begin{theorem}[有限个事件的加法定理]
\label{thm:study-inclusion-exclusion}
设 $n\ge1$ 为整数，$A_1,\ldots,A_n$ 是同一概率空间中的任意事件，则
\begin{equation}
  \mathbb P\!\left(\bigcup_{i=1}^{n}A_i\right)
    =\sum_{k=1}^{n}(-1)^{k-1}
      \sum_{1\le i_1<\cdots<i_k\le n}
      \mathbb P(A_{i_1}\cap\cdots\cap A_{i_k}).
  \label{eq:study-inclusion-exclusion-general}
\end{equation}
内层求和遍历所有由 $k$ 个不同下标组成的组合，每组只取一次，共有 $\binom nk$ 项。因此，右端依次为单事件概率之和，减去两两交集概率之和，加上三者交集概率之和，直到符号为 $(-1)^{n-1}$ 的全体事件交集项。公式不要求各事件互斥，也不要求相互独立。
\end{theorem}

\begin{proof}
用数学归纳法证明。$n=1$ 时两端均为 $\mathbb P(A_1)$；$n=2$ 时就是已经证明的式~\eqref{eq:study-inclusion-exclusion-two}。设 $n\ge3$，假设公式对任意 $n-1$ 个事件成立，令
\[
  U=\bigcup_{i=1}^{n-1}A_i,\qquad B_i=A_i\cap A_n\quad(1\le i\le n-1).
\]
由分配律，$U\cap A_n=\bigcup_{i=1}^{n-1}B_i$。再用二事件公式，得
\begin{equation}
  \mathbb P\!\left(\bigcup_{i=1}^{n}A_i\right)
    =\mathbb P(U)+\mathbb P(A_n)
      -\mathbb P\!\left(\bigcup_{i=1}^{n-1}B_i\right).
  \label{eq:study-inclusion-exclusion-induction}
\end{equation}
分别对 $A_1,\ldots,A_{n-1}$ 和 $B_1,\ldots,B_{n-1}$ 应用归纳假设。$\mathbb P(U)$ 的展开给出所有不含下标 $n$ 的交集项，$\mathbb P(A_n)$ 给出只含下标 $n$ 的一阶项。最后一项中，每个 $k$ 重交集满足
\[
  B_{i_1}\cap\cdots\cap B_{i_k}
    =A_{i_1}\cap\cdots\cap A_{i_k}\cap A_n,
  \qquad 1\le k\le n-1.
\]
它对应一个含下标 $n$ 的 $k+1$ 重交集。归纳公式中的系数为 $(-1)^{k-1}$，再乘以式~\eqref{eq:study-inclusion-exclusion-induction} 中的负号，变为 $(-1)^k$，恰好是 $k+1$ 重交集应有的符号。所有非空下标组合由此不重不漏地出现，合并便得到式~\eqref{eq:study-inclusion-exclusion-general}，归纳完成。
\end{proof}

容斥原理与对立事件法可以结合使用。由德摩根律~\eqref{eq:study-de-morgan-finite}，“所有事件都不发生”是“至少一个事件发生”的对立事件，因而
\begin{equation}
  \begin{aligned}
    \mathbb P\!\left(\bigcap_{i=1}^{n}A_i^{\mathrm c}\right)
      &=1-\mathbb P\!\left(\bigcup_{i=1}^{n}A_i\right)\\
      &=1-\sum_{k=1}^{n}(-1)^{k-1}
        \sum_{1\le i_1<\cdots<i_k\le n}
        \mathbb P(A_{i_1}\cap\cdots\cap A_{i_k}).
  \end{aligned}
  \label{eq:study-inclusion-exclusion-complement}
\end{equation}

\begin{example}
\label{ex:study-inclusion-exclusion-complement}
设 $A,B,C$ 是同一概率空间中的三个事件，已知
\[
  \mathbb P(A)=\mathbb P(B)=\frac14,\qquad \mathbb P(C)=\frac25,
  \qquad \mathbb P(A\cap B)=0,\qquad
  \mathbb P(A\cap C)=\mathbb P(B\cap C)=\frac16.
\]
求三个事件都不发生的概率 $\mathbb P(A^{\mathrm c}\cap B^{\mathrm c}\cap C^{\mathrm c})$。
\end{example}

\begin{solve}
先确定题目没有直接给出的三者交集概率。由于 $A\cap B\cap C\subseteq A\cap B$，由概率的单调性，
\[
  0\le\mathbb P(A\cap B\cap C)\le\mathbb P(A\cap B)=0,
\]
所以 $\mathbb P(A\cap B\cap C)=0$。结合对立事件关系和三事件公式~\eqref{eq:study-inclusion-exclusion-three}，得到
\begin{equation}
  \begin{aligned}
    \mathbb P(A^{\mathrm c}\cap B^{\mathrm c}\cap C^{\mathrm c})
      &=1-\mathbb P(A\cup B\cup C)\\
      &=1-\left[\frac14+\frac14+\frac25-0-\frac16-\frac16+0\right]\\
      &=1-\frac{17}{30}=\boxed{\frac{13}{30}}.
  \end{aligned}
  \label{eq:study-inclusion-exclusion-example-answer}
\end{equation}
条件 $\mathbb P(A\cap B)=0$ 不能在一般概率空间中推出 $A\cap B=\varnothing$，所以不宜直接把 $A,B$ 称为互斥事件；它已经足以推出三者交集的概率为零。此外，本题中 $\mathbb P(A)\mathbb P(B)=1/16\ne\mathbb P(A\cap B)$，故 $A,B$ 并不独立，也不能用三个对立事件概率的乘积代替所求概率。
\end{solve}

\section{条件概率}
\label{sec:study-conditional-probability}

当已知事件 $B$ 发生时，计算事件 $A$ 发生的概率，需要把考虑范围限制在 $B$ 内，并以 $B$ 的概率作为新的归一化基准。在这一范围内，$A$ 发生对应的部分是 $A\cap B$，因此应比较 $A\cap B$ 的概率与 $B$ 的概率。

\begin{definition}
\label{def:study-conditional-probability}
设 $A,B$ 是同一概率空间中的两个事件，且 $\mathbb P(B)>0$。称
\begin{equation}
  \mathbb P(A\mid B)
    =\frac{\mathbb P(A\cap B)}{\mathbb P(B)}
  \label{eq:study-conditional-probability}
\end{equation}
为在事件 $B$ 发生的条件下事件 $A$ 发生的\keyterm{条件概率（conditional probability）}。竖线右侧的 $B$ 表示已知条件，左侧的 $A$ 表示所考察的事件。
\end{definition}

定义中的分母用于将条件事件 $B$ 的概率归一化为 $1$，即 $\mathbb P(B\mid B)=1$。又因为 $A\cap B\subseteq B$，所以 $0\le\mathbb P(A\cap B)\le\mathbb P(B)$，除以正数 $\mathbb P(B)$ 得 $0\le\mathbb P(A\mid B)\le1$。特别地，在有限等可能的古典概型中，若 $B\ne\varnothing$，由式~\eqref{eq:study-classical-probability} 可得
\[
  \mathbb P(A\mid B)
    =\frac{|A\cap B|/|\Omega|}{|B|/|\Omega|}
    =\frac{|A\cap B|}{|B|},
\]
其中 $\Omega$ 为原样本空间，$|\cdot|$ 表示集合中元素的个数。这说明在古典概型中，可以只在 $B$ 包含的样本点中重新计数；一般概率模型则应使用概率之比，不能直接以样本点数之比代替。

条件 $\mathbb P(B)>0$ 不可省略。当 $\mathbb P(B)=0$ 时，由单调性也有 $\mathbb P(A\cap B)=0$，式~\eqref{eq:study-conditional-probability} 的右端成为没有定义的 $0/0$，因此本定义不能用于概率为零的条件事件。这并不意味着 $B$ 必为空集。还应注意，$\mathbb P(A\mid B)$ 与 $\mathbb P(A\cap B)$ 含义不同：前者是在已知 $B$ 发生后计算 $A$ 的概率，后者是在原概率模型中计算两者同时发生的概率。当 $\mathbb P(A)>0$ 时，$\mathbb P(B\mid A)$ 则交换了条件与所考察事件，通常也不等于 $\mathbb P(A\mid B)$；条件关系本身不表示时间先后或因果关系。

\begin{example}
\label{ex:study-conditional-probability-die}
掷一枚均匀骰子，令 $A$ 表示“点数大于 $3$”，$B$ 表示“点数为偶数”。已知出现偶数，求点数大于 $3$ 的条件概率。
\end{example}

\begin{solve}
样本空间为 $\Omega=\{1,2,3,4,5,6\}$，各样本点等可能。由于
\[
  A=\{4,5,6\},\qquad B=\{2,4,6\},\qquad A\cap B=\{4,6\},
\]
有 $\mathbb P(B)=3/6>0$、$\mathbb P(A\cap B)=2/6$，故由定义得
\[
  \mathbb P(A\mid B)=\frac{2/6}{3/6}=\frac23.
\]
也可以直接在条件事件 $B$ 的三个等可能样本点 $2,4,6$ 中计数，其中 $4,6$ 满足 $A$，所以结果为 $2/3$。原来的无条件概率为 $\mathbb P(A)=3/6=1/2$，而已知 $B$ 发生后的条件概率为 $2/3$，说明获得条件信息后，同一事件的概率可能改变。
\end{solve}

\section{独立性}
\label{sec:study-independence}
% 待补充本节的学习思考与习题解答。

\section{扩展进阶：求概率的一些方法}
\label{sec:study-probability-methods}
% 教材将本节标注为带星号的扩展进阶内容。
% 待补充本节的学习思考与习题解答。

\section{扩展阅读1：贝叶斯公式和垃圾邮件识别}
\label{sec:study-bayes-spam}
% 待补充本节的学习思考与习题解答。

\section{扩展阅读2：三门问题}
\label{sec:study-monty-hall}
% 待补充本节的学习思考与习题解答。

\section*{本章总结}
\addcontentsline{toc}{section}{本章总结}
% 待补充；整理习题时注明教材题号。

\section*{习题}
\addcontentsline{toc}{section}{习题}
% 待补充；整理习题时注明教材题号。

\chapter{随机变量及其分布}
\label{chap:study-random-variables}

\section{随机变量的概念}
\label{sec:study-random-variable-concept}
% 待补充本节的学习思考与习题解答。

\section{离散型随机变量的分布}
\label{sec:study-discrete-distributions}
% 待补充本节的学习思考与习题解答。

\section{连续型随机变量的分布}
\label{sec:study-continuous-distributions}
% 待补充本节的学习思考与习题解答。

\section{随机变量函数的分布}
\label{sec:study-transformed-distributions}
% 待补充本节的学习思考与习题解答。

\section{扩展阅读：正态分布的由来}
\label{sec:study-normal-origin}
% 待补充本节的学习思考与习题解答。

\section*{本章总结}
\addcontentsline{toc}{section}{本章总结}
% 待补充；整理习题时注明教材题号。

\section*{习题}
\addcontentsline{toc}{section}{习题}
% 待补充；整理习题时注明教材题号。

\chapter{多维随机变量及其分布}
\label{chap:study-multivariate}

\section{多维随机变量及其分布}
\label{sec:study-joint-distributions}
% 待补充本节的学习思考与习题解答。

\section{边缘（际）分布}
\label{sec:study-marginal-distributions}
% 待补充本节的学习思考与习题解答。

\section{条件分布}
\label{sec:study-conditional-distributions}
% 待补充本节的学习思考与习题解答。

\section{相互独立的随机变量}
\label{sec:study-independent-variables}
% 待补充本节的学习思考与习题解答。

\section{随机向量函数的分布}
\label{sec:study-vector-transformations}
% 待补充本节的学习思考与习题解答。

\section{扩展阅读：辛普森悖论}
\label{sec:study-simpson-paradox}
% 待补充本节的学习思考与习题解答。

\section*{本章总结}
\addcontentsline{toc}{section}{本章总结}
% 待补充；整理习题时注明教材题号。

\section*{习题}
\addcontentsline{toc}{section}{习题}
% 待补充；整理习题时注明教材题号。

\chapter{随机变量的数字特征和极限定理}
\label{chap:study-moments-limits}

\section{数学期望和中位数}
\label{sec:study-expectation-median}
% 待补充本节的学习思考与习题解答。

\section{方差和矩}
\label{sec:study-variance-moments}
% 待补充本节的学习思考与习题解答。

\section{熵的基本概念}
\label{sec:study-entropy}
% 待补充本节的学习思考与习题解答。

\section{大数定律和中心极限定理}
\label{sec:study-limit-theorems}
% 待补充本节的学习思考与习题解答。

\section{扩展阅读：数学期望的计算}
\label{sec:study-expectation-calculation}
% 待补充本节的学习思考与习题解答。

\section*{本章总结}
\addcontentsline{toc}{section}{本章总结}
% 待补充；整理习题时注明教材题号。

\section*{习题}
\addcontentsline{toc}{section}{习题}
% 待补充；整理习题时注明教材题号。

\chapter{统计学基本概念}
\label{chap:study-statistics-basics}

\section{统计学发展简史}
\label{sec:study-statistics-history}
% 待补充本节的学习思考与习题解答。

\section{基本概念}
\label{sec:study-statistical-concepts}
% 待补充本节的学习思考与习题解答。

\section{抽样分布}
\label{sec:study-sampling-distributions}
% 待补充本节的学习思考与习题解答。

\section{扩展阅读1：民意调查}
\label{sec:study-opinion-polls}
% 待补充本节的学习思考与习题解答。

\section{扩展阅读2：双盲对照试验}
\label{sec:study-double-blind-trials}
% 待补充本节的学习思考与习题解答。

\section*{本章总结}
\addcontentsline{toc}{section}{本章总结}
% 待补充；整理习题时注明教材题号。

\section*{习题}
\addcontentsline{toc}{section}{习题}
% 待补充；整理习题时注明教材题号。

\chapter{参数点估计}
\label{chap:study-point-estimation}

\section{参数点估计的概念}
\label{sec:study-point-estimation-concept}
% 待补充本节的学习思考与习题解答。

\section{矩估计法}
\label{sec:study-method-of-moments}
% 待补充本节的学习思考与习题解答。

\section{最大似然估计}
\label{sec:study-maximum-likelihood}
% 待补充本节的学习思考与习题解答。

\section{优良性准则}
\label{sec:study-estimator-criteria}
% 待补充本节的学习思考与习题解答。

\section{点估计量的大样本理论}
\label{sec:study-estimator-asymptotics}
% 待补充本节的学习思考与习题解答。

\section{扩展阅读：德军坦克问题}
\label{sec:study-german-tank}
% 待补充本节的学习思考与习题解答。

\section*{本章总结}
\addcontentsline{toc}{section}{本章总结}
% 待补充；整理习题时注明教材题号。

\section*{习题}
\addcontentsline{toc}{section}{习题}
% 待补充；整理习题时注明教材题号。

\chapter{区间估计}
\label{chap:study-interval-estimation}

\section{基本概念}
\label{sec:study-interval-concepts}
% 待补充本节的学习思考与习题解答。

\section{枢轴变量法}
\label{sec:study-pivotal-quantity}
% 待补充本节的学习思考与习题解答。

\section{大样本方法}
\label{sec:study-large-sample-intervals}
% 待补充本节的学习思考与习题解答。

\section{自助法置信区间}
\label{sec:study-bootstrap-intervals}
% 待补充本节的学习思考与习题解答。

\section{置信限}
\label{sec:study-confidence-bounds}
% 待补充本节的学习思考与习题解答。

\section{扩展阅读：“足球赛会杀人”的真假}
\label{sec:study-football-claim}
% 待补充本节的学习思考与习题解答。

\section*{本章总结}
\addcontentsline{toc}{section}{本章总结}
% 待补充；整理习题时注明教材题号。

\section*{习题}
\addcontentsline{toc}{section}{习题}
% 待补充；整理习题时注明教材题号。

\chapter{假设检验}
\label{chap:study-hypothesis-testing}

\section{问题的提法和基本概念}
\label{sec:study-testing-concepts}
% 待补充本节的学习思考与习题解答。

\section{正态总体参数检验}
\label{sec:study-normal-tests}
% 待补充本节的学习思考与习题解答。

\section{比例 $p$ 的检验}
\label{sec:study-proportion-test}
% 待补充本节的学习思考与习题解答。

\section{似然比检验}
\label{sec:study-likelihood-ratio}
% 待补充本节的学习思考与习题解答。

\section{$p$ 值}
\label{sec:study-p-value}
% 待补充本节的学习思考与习题解答。

\section{扩展阅读：多重假设检验}
\label{sec:study-multiple-testing}
% 待补充本节的学习思考与习题解答。

\section*{本章总结}
\addcontentsline{toc}{section}{本章总结}
% 待补充；整理习题时注明教材题号。

\section*{习题}
\addcontentsline{toc}{section}{习题}
% 待补充；整理习题时注明教材题号。

\chapter{非参数假设检验}
\label{chap:study-nonparametric-testing}

\section{拟合优度检验}
\label{sec:study-goodness-of-fit}
% 待补充本节的学习思考与习题解答。

\section{威尔科克森秩和检验}
\label{sec:study-wilcoxon}
% 待补充本节的学习思考与习题解答。

\section{符号检验}
\label{sec:study-sign-test}
% 待补充本节的学习思考与习题解答。

\section{其他非参数检验概述}
\label{sec:study-other-nonparametric}
% 待补充本节的学习思考与习题解答。

\section{扩展阅读：正态性检验}
\label{sec:study-normality-tests}
% 待补充本节的学习思考与习题解答。

\section*{本章总结}
\addcontentsline{toc}{section}{本章总结}
% 待补充；整理习题时注明教材题号。

\section*{习题}
\addcontentsline{toc}{section}{习题}
% 待补充；整理习题时注明教材题号。

\chapter{相关分析和回归分析}
\label{chap:study-correlation-regression}

\section{相关分析}
\label{sec:study-correlation}
% 待补充本节的学习思考与习题解答。

\section{回归分析}
\label{sec:study-regression}
% 待补充本节的学习思考与习题解答。

\section{多元回归中自变量的选择和模型诊断简述}
\label{sec:study-regression-diagnostics}
% 待补充本节的学习思考与习题解答。

\section{扩展阅读：相关与因果}
\label{sec:study-correlation-causation}
% 待补充本节的学习思考与习题解答。

\section{附录}
\label{sec:study-regression-appendix}
% 待补充本节的学习思考与习题解答。

\section*{本章总结}
\addcontentsline{toc}{section}{本章总结}
% 待补充；整理习题时注明教材题号。

\section*{习题}
\addcontentsline{toc}{section}{习题}
% 待补充；整理习题时注明教材题号。

\end{document}
